\section{Conclusion}

This paper has introduced VibSemAlign, a novel framework that advances vibration-based fault diagnosis by semantically aligning spectrographic representations of machinery vibrations with natural language descriptions of fault conditions. By leveraging vision-language models fine-tuned through a specialized vibration-semantic contrastive loss and integrated with model-agnostic meta-learning, VibSemAlign achieves robust zero-shot and few-shot performance across diverse domains, as demonstrated on benchmarks such as CWRU and Paderborn datasets. Our approach addresses longstanding challenges in predictive maintenance, including data scarcity, domain shifts due to varying operational conditions, and the need for interpretable diagnostics.

The core contributions of this work are threefold. First, we propose the vibration-semantic contrastive loss (VSCL), which effectively bridges visual signal patterns with textual semantics, enabling the model to discern subtle fault signatures like bearing defect modulations without extensive retraining. Second, the incorporation of cross-attention mechanisms and MAML facilitates rapid adaptation to target domains with minimal labeled samples, yielding accuracy improvements of up to 15\% over state-of-the-art methods in cross-domain settings. Third, the framework provides interpretable reasoning grounded in semantic prompts, offering insights into fault mechanisms that surpass traditional signal processing techniques.

Empirical validation on CWRU and Paderborn benchmarks confirms VibSemAlign's superiority: 96.8\% accuracy on CWRU (surpassing CNN-1D by 7.6\% and Transformer by 5.3\%) and 94.5\% on Paderborn (improving 8--12\% over baselines). Zero-shot cross-domain transfer from CWRU to Paderborn achieves 85.3\% accuracy, outperforming DANN by 17.2 percentage points without any target domain exposure. Ablation studies isolate the contributions of each component: the contrastive loss contributes 3.6\% accuracy gain, MAML adds 2.7\%, and cross-attention fusion provides an additional 1.7\%, with synergistic effects yielding the full model's SOTA performance. These results validate the hypothesis that physics-informed semantic alignment substantially enhances generalization beyond what unimodal or generic VLM approaches can achieve.

Looking ahead, several avenues warrant further exploration. Extending VibSemAlign to multimodal fusion—incorporating acoustic emissions, motor current signatures, or thermal imaging—could enhance diagnostic granularity for complex machinery exhibiting compound faults. Real-time deployment on edge devices remains a priority, necessitating model distillation and INT8 quantization to reduce the current 52 ms inference latency to sub-10 ms targets for high-speed rotating machinery monitoring. Additionally, scaling to larger, real-world industrial datasets spanning diverse machine types (pumps, compressors, wind turbines) and exploring unsupervised domain adaptation would broaden applicability. Auto-prompting strategies that generate physics-informed descriptions from signal statistics could further reduce reliance on expert knowledge, democratizing deployment across industries lacking specialized vibration engineers.

In broader terms, VibSemAlign holds significant potential for industrial applications, enabling proactive fault detection that minimizes downtime, reduces maintenance costs, and enhances equipment reliability in sectors like manufacturing, energy production, and transportation. By bridging the gap between signal processing expertise and modern language-vision AI, this work opens a new paradigm for semantically grounded predictive maintenance systems capable of generalizing across the diverse and challenging conditions of real-world industrial environments.

